\clearpage
\VerdierPhysicalPageAnchor{0200}%
\begingroup\small\setlength{\parskip}{0pt}\setlength{\emergencystretch}{1.5em}
\setlength{\abovedisplayskip}{4pt plus 2pt minus 2pt}%
\setlength{\belowdisplayskip}{4pt plus 2pt minus 2pt}%

\VerdierSemanticAnchor{ch3d-2-3-3-condition-2}
\noindent\textit{2) The sequence $i\longmapsto i-M\text{-}\dim(Y^i)$ tends
to $+\infty$ as $i$ tends to $+\infty$.}

\VerdierSemanticAnchor{ch3d-2-3-3-quasi-isomorphism}
\textit{Then there is a complex
$X\in\operatorname{Ob}\operatorname{comp}(\mathcal A)$ whose components
belong to $M$, and a quasi-isomorphism}
\[
u\colon X\longrightarrow Y.
\]

\VerdierSemanticAnchor{ch3d-2-3-3-bounded-cases}
\textit{If, in addition, all the components of $Y$ have finite
$M$-dimension and
$Y\in\operatorname{Ob}(\operatorname{comp}^{b}(\mathcal A))$
(respectively, $\operatorname{Ob}(\operatorname{comp}^{+}(\mathcal A))$),
then $X$ may be taken in $\operatorname{comp}^{b}(\mathcal A)$
(respectively, $\operatorname{comp}^{+}(\mathcal A)$).}

\VerdierSemanticAnchor{ch3d-2-3-3-proof}
Let $n_0$ be an integer such that
$i>n_0\Rightarrow M\text{-}\dim(Y^i)<+\infty$. We propose to construct:

\VerdierSemanticAnchor{ch3d-2-3-3-induction-data}
\noindent a) A sequence of complexes $X_i$, for $i\geq n_0$.

\noindent b) A sequence of morphisms
$u_i\colon X_i\longrightarrow X_{i-1}$ for $i>n_0$, and
$u_{n_0}\colon X_{n_0}\longrightarrow Y$, such that:

\VerdierSemanticAnchor{ch3d-2-3-3-properties}
\noindent $1)_i$ $(X_i)^j\in M$ for $j\leq i$.

\noindent $2)_i$ $(X_i)^j=Y^j$ for $j>i$.

\noindent $3)_i$ The morphism $u_i$ is a quasi-isomorphism.

\noindent $4)_i$ $(u_i)^j=\operatorname{id}$ for $j>i$.

\noindent $5)_i$ $(u_i)^j=\operatorname{id}$ for
$j\leq i-M\text{-}\dim(Y^i)$; this last property is asserted only for
$i>n_0$.

\VerdierSemanticAnchor{ch3d-2-3-3-projective-limit}
Once this construction has been made, observe that the inverse system $X_i$
has an inverse limit in $\operatorname{comp}(\mathcal A)$: the inverse
systems induced on the components of any fixed degree are stationary, since
$i-M\text{-}\dim(Y^i)$ tends to $+\infty$ with $i$. It will therefore suffice
to set $X=\varprojlim_i X_i$ and take as
$u\colon X\longrightarrow Y$ the canonical morphism. The components of the
complex $X$ belong to $M$, and $u$ is a quasi-isomorphism.

\VerdierSemanticAnchor{ch3d-2-3-3-un0}
First construct the morphism $u_{n_0}\colon X_{n_0}\longrightarrow Y$.
Proceed inductively, and suppose that for every integer $i$ with
$-n_0\leq i<q$, complexes $X_{n_0,i}$ and morphisms
$v_{n_0,i}\colon X_{n_0,i}\longrightarrow X_{n_0,i-1}$ for $i>-n_0$, and
$v_{n_0,-n_0}\colon X_{n_0,-n_0}\longrightarrow Y$, have been constructed
such that:

\VerdierSemanticAnchor{ch3d-2-3-3-inner-properties}
\noindent $I)_{n_0,i}$ The components $(X_{n_0,i})^j$ belong to $M$ for
$-i\leq j\leq n_0$.

\noindent $II)_{n_0,i}$ The components $(X_{n_0,i})^j$ belong to $M'$ for
every integer $j$.

\noindent $III)_{n_0,i}$ The morphism $v_{n_0,i}$ is an elementary morphism
of height $-i$ (2.1.3).

\VerdierSemanticAnchor{ch3d-2-3-3-xn0qminus1-shape}
The complex $X_{n_0,q-1}$ then has the form
\[
\resizebox{\textwidth}{!}{$\displaystyle
\cdots\longrightarrow Y^{-q-1}\longrightarrow Z^{-q}
\longrightarrow L^{-q+1}\longrightarrow L^{-q+2}\longrightarrow\cdots
\longrightarrow L^{n_0-1}\longrightarrow L^{n_0}
\longrightarrow Y^{n_0+1}\longrightarrow Y^{n_0+2}\longrightarrow\cdots,
$}
\]
\endgroup
